\section{Method}
\label{sec:method}

To isolate how the pruning criterion affects post-RLVR recovery, we hold everything else constant: a shared structural analysis computes importance and functional overlap scores for all prunable structures, four scoring criteria combine these scores differently to select structures for removal, and an identical two-stage recovery pipeline is applied to every pruned model.
The only variable across conditions is which structures are removed (Figure~\ref{fig:method_framework}).

\subsection{Preliminaries and Notation}
\label{sec:prelim}

\paragraph{Structured pruning.}
Given a pretrained transformer $\mathcal{M}$ with $L$ layers, structured pruning removes entire computational units (attention heads or MLP channels) to reduce model size while maintaining hardware-friendly dense computation.
We denote the set of prunable structures as $\mathcal{H} = \{h_1, h_2, \ldots, h_N\}$.
In this work, we focus exclusively on MLP channel group pruning; the formalism below covers attention heads for completeness, but all experiments remove only MLP groups.

\paragraph{Grouped Query Attention (GQA).}
In GQA architectures~\citep{ainslie2023gqa}, $n_Q$ query heads share $n_{KV}$ key-value heads ($G = n_Q / n_{KV}$ per group).
For DeepSeek-R1-Distill-Qwen-7B ($n_Q{=}28$, $n_{KV}{=}4$, $G{=}7$), we treat each GQA group (7 Q heads + 1 KV head) as a single prunable unit.

\paragraph{RLVR and GRPO.}
Reinforcement Learning with Verifiable Rewards (RLVR) recovers pruned model capabilities using outcome-based rewards from tasks with verifiable answers.
We use Group Relative Policy Optimization (GRPO;~\citealp{shao2024grpo}) as the recovery algorithm.

\subsection{Structural Analysis}
\label{sec:structural_analysis}

All four pruning criteria share a common structural analysis stage that computes two scores per structure: \emph{importance} $I(h)$ capturing the structure's contribution to reasoning, and \emph{functional overlap} $F(h)$ measuring representational redundancy with other structures in the same layer.

\paragraph{Calibration data.}
We compute importance and overlap scores on chain-of-thought reasoning data from GSM8K~\citep{cobbe2021gsm8k} training examples.
Given a calibration set $\mathcal{D}_\text{cal} = \{(x_j, y_j)\}_{j=1}^{M}$ with $M{=}512$ randomly sampled examples, we compute activations and gradients on the concatenated sequences $(x_j, y_j)$.

\paragraph{Structural importance.}
For each structure $h$ with activation $\mathbf{a}_h$, we approximate the loss change from removal via first-order Taylor expansion ($\Delta\mathcal{L} \approx -\mathbf{a}_h^\top \nabla_{\mathbf{a}_h}\mathcal{L}$, dropping $O(\|\mathbf{a}_h\|^2)$ terms) and aggregate over the calibration set:
\begin{equation}
    I(h) = \left| \sum_{(x,y) \in \mathcal{D}_\text{cal}} \mathbf{a}_h(x, y)^\top \frac{\partial \mathcal{L}(x, y)}{\partial \mathbf{a}_h} \right|,
    \label{eq:importance}
\end{equation}
where $\mathcal{L}$ is the language modeling loss over the reasoning trace.
For a GQA group $g$ containing Q heads $\{q_1, \ldots, q_G\}$ and KV head $k$, we aggregate: $I(g) = \sum_{i=1}^{G} I(q_i) + I(k)$.

\paragraph{Functional overlap via CKA.}
We measure functional overlap using linear Centered Kernel Alignment (CKA;~\citealp{kornblith2019cka}), which quantifies similarity between two sets of representations invariant to orthogonal transformations and isotropic scaling.
Given activation matrices $\mathbf{X}, \mathbf{Y} \in \mathbb{R}^{n \times d}$ collected over $n$ calibration examples, linear CKA is:
\begin{equation}
    \text{CKA}(\mathbf{X}, \mathbf{Y}) = \frac{\text{HSIC}(\mathbf{X}, \mathbf{Y})}{\sqrt{\text{HSIC}(\mathbf{X}, \mathbf{X}) \cdot \text{HSIC}(\mathbf{Y}, \mathbf{Y})}},
    \label{eq:cka}
\end{equation}
where HSIC is the Hilbert-Schmidt Independence Criterion, $\text{HSIC}(\mathbf{X}, \mathbf{Y}) = \frac{1}{(n-1)^2} \operatorname{tr}(\mathbf{K} \mathbf{H} \mathbf{L} \mathbf{H})$, with $\mathbf{K} = \mathbf{X}\mathbf{X}^\top$, $\mathbf{L} = \mathbf{Y}\mathbf{Y}^\top$, and $\mathbf{H} = \mathbf{I}_n - \frac{1}{n}\mathbf{1}\mathbf{1}^\top$.

For a GQA group $g_i$ in layer $\ell$, we form its representation by concatenating its constituent Q head activations: $\mathbf{X}_{g_i} = [\mathbf{a}_{q_1}; \ldots; \mathbf{a}_{q_G}] \in \mathbb{R}^{n \times Gd_h}$.
The functional overlap of $g_i$ is the mean pairwise CKA with other groups in the same layer:
\begin{equation}
    F(g_i) = \frac{1}{|\mathcal{G}_\ell| - 1} \sum_{\substack{g_j \in \mathcal{G}_\ell \\ j \neq i}} \text{CKA}(\mathbf{X}_{g_i}, \mathbf{X}_{g_j}).
    \label{eq:overlap}
\end{equation}
We apply the same intra-layer CKA overlap to MLP channel groups, partitioning each MLP intermediate dimension into equal-sized groups.
Unlike prior work that uses inter-layer CKA to identify globally redundant layers for depth pruning~\citep{mpruner2024,lachance2022cka}, we use \emph{intra-layer} CKA between individual structures to measure per-structure functional redundancy.
Computing all pairwise CKA scores adds minimal overhead ($O(L \cdot k^2)$, $k \le 28$; $<$2 min on one GPU).

\paragraph{CKA reliability.}
Our intra-layer use between structures of identical dimensionality avoids sensitivities identified for cross-layer comparisons~\citep{davari2023cka,klabunde_similarity_survey_2023}; rankings are stable across calibration sizes (Spearman $\rho \geq 0.964$; Appendix~\ref{sec:appendix_cka_stability}).

\subsection{Pruning Scoring Criteria}
\label{sec:scoring}

We compare four scoring criteria that combine importance $I(h)$ and functional overlap $F(h)$ in different ways.
All four methods use the same per-layer capacity floor (${\ge}$40\% of MLP groups per layer must be retained, chosen to prevent any single layer from becoming a throughput bottleneck; ablation over $\tau \in \{0.3, 0.4, 0.5, 0.6\}$ in Appendix~\ref{sec:appendix_floor}) and global ranking procedure: structures are sorted by score, and the lowest-scoring structures are pruned up to the target sparsity ratio, skipping any candidate whose removal would violate the layer floor.

\paragraph{Standard (importance-only).}
Following LLM-Pruner~\citep{llmpruner2023}, structures are ranked solely by importance:
\begin{equation}
    S_\text{std}(h) = I(h).
    \label{eq:standard}
\end{equation}
This is the industry-standard criterion that retains gradient-important structures and removes ``unimportant'' ones.

\paragraph{RASP (recovery-aware composite).}
\rasp{} is a parameterized family of scoring functions indexed by $\alpha$, which controls the trade-off between gradient importance and diversity preservation:
\begin{equation}
    S_\text{rasp}(h) = I(h) \times \big(1 - F(h)\big)^\alpha,
    \label{eq:score}
\end{equation}
where $\alpha{=}0$ reduces to Standard (pure importance) and larger $\alpha$ increasingly favors diversity.
Dose-response analysis (\S\ref{sec:sparsity_cliff}) reveals that $\alpha{=}1.0$ yields suboptimal results---exhibiting reward collapse under GRPO due to insufficient diversity weighting---while $\alpha{=}2.0$ achieves peak recovery.
The multiplicative form ensures that structures with zero importance are automatically pruned ($S{=}0$) regardless of overlap; an additive alternative $\beta I(h) + (1{-}\beta)(1{-}F(h))$ would retain zero-importance structures if their overlap is low.
Structures with high importance \emph{and} low overlap receive the highest scores and are retained; those with high overlap are penalized and pruned preferentially.
At the empirically optimal $\alpha{=}2.0$, the diversity term $(1{-}F(h))^2$ dominates the score, making \rasp{} at $\alpha{=}2.0$ effectively CKA-guided with importance as a secondary filter (\S\ref{sec:sparsity_cliff}).
Nevertheless, \rasp{} at $\alpha{=}2.0$ outperforms CKA-only under both recovery paradigms (GRPO: 66.11\% vs.\ 64.39\%, $+$1.72\,pp; SFT: 64.97\% vs.\ 63.91\%, $+$1.06\,pp; Table~\ref{tab:alpha}), indicating that importance filtering provides complementary signal beyond diversity preservation alone.


\paragraph{CKA-only (overlap-only).}
CKA-only scoring disregards gradient importance, ranking structures by functional overlap alone.
Structures with the highest mean pairwise CKA similarity to other structures in their layer are removed first:
\begin{equation}
    S_\text{cka}(h) = 1 - F(h).
    \label{eq:cka_only}
\end{equation}
This criterion preserves all representationally unique structures regardless of their gradient importance, removing only those whose function is redundant with retained structures.
$S_\text{cka}$ corresponds to the limit $\alpha \to \infty$ in Eq.~\ref{eq:score}.

\paragraph{Random.}
Structures are selected uniformly at random for removal, subject to the per-layer capacity floor.
This baseline controls for the effect of any informed scoring criterion.

The four criteria span a spectrum from pure gradient importance (Standard) to pure representational diversity preservation (CKA-only), with \rasp{} at $\alpha{=}2.0$ as a diversity-dominant interpolation and Random as an uninformed control.
The key distinction is that CKA-only retains \emph{functionally unique} structures regardless of gradient contribution: structures deemed ``unimportant'' by gradient metrics may provide non-redundant computational pathways that facilitate RL exploration~\citep{rlsqueezes2026}.

\subsection{Two-Stage Recovery}
\label{sec:recovery}

After pruning, all four models undergo identical two-stage recovery: (1)~a brief SFT warmup on chain-of-thought demonstrations restores end-of-sequence generation and format compliance (pruning disrupts termination behavior, making direct GRPO infeasible); (2)~GRPO with binary reward enables RL-based exploration, redistributing functionality across remaining structures~\citep{rlsqueezes2026}, in contrast to SFT which imitates fixed trajectories.
Details in Section~\ref{sec:exp_setup} and Appendix~\ref{sec:appendix_training}.

\subsection{Recovery Capacity Diagnostics}
\label{sec:diagnostics}

We define \textbf{recovery capacity} as the model's ability to benefit from post-pruning recovery training, operationalized as peak mean reward over the GRPO training horizon.
Four diagnostics (mean reward, zero-variance batch fraction ($f_0$), gradient norm, and Kullback--Leibler (KL) divergence) distinguish three regimes: \emph{active learning} (rising reward, low $f_0$), \emph{stagnation} (flat reward), and \emph{collapse} (declining reward, high $f_0$, vanishing gradients).

The complete pruning comparison procedure is given in Algorithm~\ref{alg:rasp} (Appendix~\ref{sec:appendix_algo}).
